\section{A uniform exclusion for large weighted means}
\label{mo:section}

The following estimates give a further restriction on a possible raw dual
witness. They apply to arbitrary clipped weights, independently of the
structure of the source set.
Let $M\ge4$, let $0\le w_p\le1$ be supported on primes $p\le M$, and set
\[
 X(a)=\sum_{p\le M}w_pv_p(a),\qquad
 \mu=\sum_{p\le M}\frac{w_p}{p-1},\qquad
 \nu=\sum_{p\le M}\frac{w_p^2p}{(p-1)^2}.
\]
Expectations in this section are over a uniform integer in $\{1,\ldots,M\}$.
Put $J=\lfloor\log_2M\rfloor$ and
\begin{equation}
 \epsilon_M=
 \frac{\pi(M)+(J+1)\pi(\sqrt M)}{M}
       +\frac1{\lfloor\sqrt M\rfloor}.
 \label{mo:error}
\end{equation}

\begin{lemma}[Uniform first and second moments]
\label{mo:moments}
One has
\begin{equation}
 \mu-\epsilon_M\le \mathbb EX\le\mu,\qquad
 \mathbb EX^2\le\mu^2+\nu,\qquad
 \nu\le\min\left\{2\mu,\frac76\mu+\frac{49}{48}\right\}.
 \label{mo:bounds}
\end{equation}
Moreover,
\[
 \epsilon_M\le\frac{\pi(M)}M+\frac{J+3}{\sqrt M}
 =O\left(\frac1{\log M}+\frac{\log M}{\sqrt M}\right)=o(1).
\]
\end{lemma}

\begin{proof}
Let $j_p=\lfloor\log M/\log p\rfloor$. Prime-power counting gives
\[
 \mathbb EX=\frac1M\sum_{p\le M}w_p
                    \sum_{j=1}^{j_p}\left\lfloor\frac M{p^j}\right\rfloor.
\]
Compare this with $\mu=\sum_pw_p\sum_{j\ge1}p^{-j}$.
The first-level floor errors total at most $\pi(M)/M$.
At levels $j\ge2$, only primes at most $\sqrt M$ occur, each with
at most $J-1$ further floor errors. For those primes the omitted tail is
\[
 \sum_{j>j_p}p^{-j}
  =\frac1{p^{j_p}(p-1)}
  <\frac p{M(p-1)}\le\frac2M.
\]
For $p>\sqrt M$, the tail starts at level two, and
\[
 \sum_{p>\sqrt M}\frac1{p(p-1)}
 \le\sum_{k>\sqrt M}\frac1{k(k-1)}
 =\frac1{\lfloor\sqrt M\rfloor}.
\]
Adding these bounds proves the mean estimate. The asserted error
asymptotic follows from Chebyshev's upper bound for $\pi$;
see, for example, \cite{GranvilleSound}.

For distinct primes $p,q$, divisibility by both levels $p^j,q^k$ gives
\[
 \mathbb E[v_p(a)v_q(a)]
 \le\sum_{j,k\ge1}\frac1{p^jq^k}
 =\frac1{(p-1)(q-1)}.
\]
At one prime,
\[
 \mathbb E[v_p(a)^2]
 \le\sum_{j\ge1}\frac{2j-1}{p^j}
 =\frac{p+1}{(p-1)^2}.
\]
Expanding $X^2$ proves $\mathbb EX^2\le\mu^2+\nu$.
No independence assumption is used.
The inequality $\nu\le2\mu$ follows from $w_p^2\le w_p$ and
$p/(p-1)\le2$.

For the sharper envelope, every $p\ge7$ contributes at most
$(7/6)w_p/(p-1)$ to $\nu$. The maximum possible excesses at $p=2,3,5$
are respectively $5/6$, $1/6$, and $1/48$. Their sum is $49/48$,
proving the last part of \eqref{mo:bounds}.
\end{proof}

\begin{theorem}[A sufficient mean threshold]
\label{mo:threshold}
For all $M$ with $\epsilon_M\le1/2$, the condition
\begin{equation}
 \mu\ge\frac{241}{136}+4\epsilon_M
 \label{mo:threshold-simple}
\end{equation}
implies, for every positive $M$-term interval $I$,
\begin{equation}
 \sum_{a=1}^M(X(a)-2)_+
 \le\sum_{b\in I}(X(b)-1)_+.
 \label{mo:h2}
\end{equation}
The limiting sufficient threshold is
$241/136=1.7720588235\ldots$.
\end{theorem}

\begin{proof}
For any nonnegative random variable with mean $m$ and second moment $q$,
\[
 \mathbb E\min(X,2)
 =\frac{m+2-\mathbb E|X-2|}{2}
 \ge\frac{m+2-\sqrt{q-4m+4}}2.
\]
Thus $q\le m^2+4m-4$ is sufficient for capped mean at least one.
Put $\epsilon=\epsilon_M$. Since $m\ge\mu-\epsilon$, it is enough that
\[
 \nu\le4\mu-4-(2\mu+4)\epsilon+\epsilon^2.
\]
The sharper envelope in Lemma~\ref{mo:moments} reduces this to
\begin{equation}
 \left(\frac{17}{6}-2\epsilon\right)\mu
 \ge\frac{241}{48}+4\epsilon-\epsilon^2.
 \label{mo:threshold-exact}
\end{equation}
At $\mu=241/136+4\epsilon$, the left side minus the right side equals
\[
 \epsilon\left(\frac{773}{204}-7\epsilon\right)\ge0
 \qquad(0\le\epsilon\le1/2).
\]
The coefficient of $\mu$ is positive, so larger $\mu$ also satisfy
\eqref{mo:threshold-exact}. In this range $\mu-\epsilon>0$, which
justifies using monotonicity of $m^2+4m$ on the nonnegative half-line.

Finally, every prime-power count in $I$ is at least its prefix count,
and hence $\sum_I X\ge\sum_{a\le M}X(a)$. Therefore
\[
 \begin{split}
 &\sum_{a\le M}(X(a)-2)_+-\sum_{b\in I}(X(b)-1)_+\\
 &\qquad\le M-\sum_{a\le M}\min(X(a),2)\le0.
 \end{split}
\]
This proves \eqref{mo:h2}.
\end{proof}

The exact sufficient threshold in \eqref{mo:threshold-exact} is
\[
 \mu\ge
 \frac{241/48+4\epsilon_M-\epsilon_M^2}{17/6-2\epsilon_M}.
\]
The simpler estimate $\mu\ge2+4\epsilon_M$ also follows directly
from $\nu\le2\mu$ and
\[
 \min(x,2)\ge ax-\frac{a^2x^2}{8}\quad(0\le a\le1),
 \qquad a=\frac4{\mu+2}.
\]
These thresholds are sufficient; optimality is not asserted.
Primes greater than $M$ must be omitted when computing $\mu$.
They have zero prefix contribution, and restoring them on the target
side only strengthens \eqref{mo:h2}.

For a distinct source subset $A\subseteq[2,M]$, the hinge inequality
implies that the corresponding raw dual objective is at most $2|A|$.
Thus a dual witness exceeding this bound must lie below the mean
threshold. This does not bound the number of substantial coordinates:
for any integer $N\ge2$, choose $N$ distinct primes greater than $N^2$
and assign each weight one. Their mean is at most $N/(N^2-1)\to0$,
despite $N$ substantial weights. This is an example of weight vectors,
not a counterexample instance. More generally, a mean bound controls only
\[
 \sum_{\{p:w_p\ge\eta\}}\frac1{p-1}\le\frac\mu\eta
 \qquad(\eta>0),
\]
and does not resolve the interaction of a growing core with the other
prime coordinates.
